\chapter{企业案例化应用说明}

\section{场景设定}
考虑某离散制造企业的两条产线：A 线以稳定大批量生产为主，B 线以小批量多规格生产为主。两条产线具有不同的刀具寿命波动与质量检验特征：
\begin{itemize}
  \item A 线：寿命均值较高、方差较小，误判率低；
  \item B 线：寿命波动较大，换型频繁，误判率相对偏高。
\end{itemize}

\section{策略迁移原则}
同一数学模型可通过参数切换适配不同产线。定义参数映射算子
\[
\pi^\star=\mathcal M(\Theta,\mathcal C),
\]
其中 $\Theta$ 为成本与统计参数，$\mathcal C$ 为产线工艺约束。迁移时应保持规则不变，仅替换参数，不直接复用旧策略节点。

\section{A线应用分析}
对 A 线，建议采用“低频检查 + 明确换刀阈值”的方案。原因是寿命分布集中，过密检查收益有限。策略目标应偏向降低检查负担并保持稳定节拍。

\section{B线应用分析}
对 B 线，建议采用“分阶段 + 双检确认”方案。前段稀疏检查保证产能，后段加密检查抑制故障扩散；对不合格判定做复检，以降低误停引起的生产中断。

\section{跨产线统一管理}
在集团层面可建立统一决策框架：
\begin{enumerate}
  \item 统一数据口径（寿命、误判、损失）；
  \item 统一求解流程（参数更新 -- 策略求解 -- A/B 测试）；
  \item 统一评估指标（单位成本、废品率、停机率）。
\end{enumerate}

\section{本章小结}
案例化分析说明本文模型可从课程作业扩展到多产线管理场景，具有较好的可迁移性与可实施性。
